\section{Set-Theoretic Types and Occurrence Typing}
\label{sec:set-theoretic}

Type slicing looks especially promising for common gradual type systems in which \textit{control flow and types are intertwined}. Here, an expression's type depends non-trivially on the branches inside a function, or the control-flow path in which the expression lies. Therefore, explaining a type genuinely requires considering the branches within the code, rather than just the basic annotations and literals. Further, as type slices necessarily \textit{retain} program structure, they also have the ability to retain appropriate control flow structure.

\subsection{Set-Theoretic Types}
\label{sec:future-set-theoretic-occurrence}

\emph{Set theoretic types} extend the type language with union, intersection, and negation. A term can then carry several types simultaneously, allowing typing of overloaded functions, and typecase constructs, commonly used in real world hybrid typed languages (e.g.\ TypeScript, or Elixir~\cite{ElixirTypes}). Relevant foundations have been explored for these types \cite{CastagnaLanvin, CastagnaLPS}, including graduality proofs \cite{AngeloThesis}, providing formal foundations to apply type slicing to.\footnote{After adapting into a bidirectional setting with expression holes.}

To demonstrate this application, consider generalising sum types of the core calculus (\S\ref{sec:core-calculus}) to allow sum injections to synthesise partially dynamic types: 
\[\inference[\synr{$\iota_1$}]{\synthesis{e}{\tau}}{\synthesis{\iota_1\; e}{\tau + \gap}} \qquad
  \inference[\synr{$\iota_2$}]{\synthesis{e}{\tau}}{\synthesis{\iota_2\; e}{\gap + \tau}}\]

Then an if statement $\mathit{pred} \triangleq \mathbf{if}\ x > 0\ \mathbf{then}\ \iota_1(x - 1)\ \mathbf{else}\ \iota_2(*)$ synthesises a sum type $\mathbb{N} + *$ without requiring type annotations in much the same way that union type inference often works (i.e. union type $\mathbb{N}\ |\ *$). Then, type slicing can intuitively explain where the components of the sum (union) come from. Querying $\gap + *$ will highlight exactly the failing branch and $\mathbb{N} + \gap$ the succeeding one. 
\subsection{Occurrence Typing}
\label{sec:occurrence}

\emph{Occurrence typing}~\cite{OccurrenceTyping} is another common technique in hybrid type systems. Inside the true branch of \texttt{if (typeof x = int) ... else ...}, the type of \texttt{x} is refined as an \texttt{int}. These explicitly require exploring the surrounding control flow, which can be extensive and/or complex, to reason about the types. As mentioned, type slices necessarily retain the program structure, so can provide a reduced slice with minimal control flow to reason about the type within. However, formally applying type slicing to such a system will need extensive work on integration with bidirectional typing and static graduality.

%%% Local Variables:
%%% mode: LaTeX
%%% TeX-master: "main"
%%% End:
